\section{The dbt Optimization Setting}
\label{sec:motivating}

Before presenting our approach, we describe the optimization setting that \tool targets and illustrate, through three examples drawn from or inspired by real-world analytics pipelines, the kinds of cross-model optimizations that this setting enables but that fall outside the reach of existing techniques.

\subsection{What is a dbt Project?}
\label{sec:dbt-setting}

A dbt project is a directed acyclic graph (DAG) of SQL transformations, called \emph{models}.
Each model is a SQL transformation that produces a single output relation;
in this paper we focus on the standard \texttt{view} and \texttt{table} materialization modes.
Models reference one another through an explicit \texttt{ref()} construct that defines
the dependency edges of the DAG.
Each model is configured to be either \emph{materialized} as a physical table or
expanded inline as a \emph{view}.
The project is executed periodically against a target data warehouse:
when a model is run, dbt resolves its \texttt{ref()} calls, substitutes the appropriate
table reference or subquery, and submits the resulting SQL to the warehouse.

This setting differs from the workloads targeted by traditional query optimization in three ways that matter for our problem.

\head{Full DAG visibility}
The optimizer has access to the complete SQL definition of every model in the project, the dependency structure that connects them, and the materialization configuration of each model.
This is a stronger form of visibility than what is available to a database query optimizer, which rewrites one query at a time, or to multi-query optimization (MQO)~\cite{sellis1988multiple}, which sees a set of independent queries.
In both cases, the optimizer reasons over a flat collection of queries; in the dbt setting, it reasons over a dependency graph of transformations whose structure is itself under its control.
This shift in scope, from query rewriting to graph-level restructuring, is what makes optimizations such as merging, splitting, or inventing new intermediate models meaningful.

\head{Project stability}
A dbt project is a maintained software artifact, not an ad-hoc query workload.
It changes only when developers modify the underlying business logic or data model, not on every execution.
Between such changes, the same DAG is executed on a recurring schedule, often many times per day.
This stability justifies investing optimization effort that would not pay off in an online query processing setting: the cost of analysis (including potentially expensive cross-model reasoning) is amortized over many scheduled executions of the same DAG.

\head{Heterogeneous refresh patterns}
Different sources update at different rates, and different downstream consumers (dashboards, reports, downstream pipelines) require fresh results at different rates.
The refresh frequency of a model is therefore not a property of the model in isolation but is shaped by both its upstream sources and its downstream consumers.
Two consequences follow.
First, a model can sometimes be skipped entirely when its upstream inputs have not changed since the last run, so the effective cost of a model depends on how often it must actually be recomputed, not on how often the schedule fires.
Second, structural decisions about the DAG are not independent of refresh scheduling: splitting a model along a frequency boundary can move expensive computation onto a slower cadence, while merging two models can force the combined result to refresh at the rate of the faster input.
Optimizing the DAG therefore requires reasoning about structure and refresh rates jointly.
 
These three properties (visibility, stability, and frequency heterogeneity) together define a distinct optimization setting.
The optimizer is reasoning not about a query but about a \emph{program}: a stable DAG of SQL transformations whose structure, materialization, and refresh schedule are jointly under its control.
A real dbt project may contain hundreds of models, and the space of possible refactorings, materialization choices, and frequency-aware splits is far too large for developers to explore systematically by hand; this motivates the need for an automated optimizer rather than a guidebook of best practices.

\subsection{Motivating Examples}
\label{sec:motivating-examples}

We illustrate three classes of optimization that \tool discovers.
Each requires the optimizer to read and reason about SQL definitions \emph{across} model boundaries within the DAG, a capability that the dbt setting uniquely affords.
The SQL in each example is simplified for presentation; domain-specific details are abstracted away.

%==============================================================
% EXAMPLE 1: Semantic Factorization
%==============================================================
\head{Example 1: Cross-Model Semantic Factorization}
A pipeline contains seven classifier models that each tag orders with a service category (\texttt{dialysis}, \texttt{psychiatric}, and five others).
The input is a \texttt{claims} table in which each row represents one item of an order, so a single order has many rows.
Each classifier scans \texttt{claims}, keeps rows whose codes match its category's code list, and returns the distinct \texttt{order\_id}s of those rows.
The seven classifiers have \emph{different} SQL: different predicates, different code vocabularies, and in some cases different join structures.
Two representative classifiers are shown on the left of Figure~\ref{fig:ex-factorization}.


\begin{figure}[!htbp]
\begin{minipage}[t]{0.48\linewidth}
\figsubheader{Before: 7 classifiers (2 shown)}
\begin{lstlisting}[language=SQL]
-- dialysis_classifier
SELECT DISTINCT order_id
FROM {{ ref('claims') }}
WHERE bill_code LIKE '72%'
   OR taxonomy IN (set_dia_tax)
   OR diagnosis IN (set_dia_dx)


-- psychiatric_classifier
SELECT DISTINCT order_id
FROM {{ ref('claims') }}
WHERE revenue IN (set_psy_rev)
   OR diagnosis BETWEEN
        '290' AND '319'
\end{lstlisting}
\vspace{1.6\baselineskip}
\end{minipage}
\hfill
\begin{minipage}[t]{0.48\linewidth}
\figsubheader{After: precomputation+lookups}
\begin{lstlisting}[language=SQL]
-- order_flags (new)
SELECT order_id,
  MAX(CASE WHEN bill_code LIKE '72%'
    OR taxonomy IN (set_dia_tax)
    OR diagnosis IN (set_dia_dx)
    THEN 1 ELSE 0 END) AS f_dia,
  MAX(CASE WHEN ... 
    THEN 1 ELSE 0 END) AS f_psy,
  ...   -- 7 categories
FROM {{ ref('claims') }}
GROUP BY order_id
 
-- dialysis_classifier (rewritten)
SELECT order_id
FROM {{ ref('order_flags') }}
WHERE f_dia = 1
\end{lstlisting}
\end{minipage}
\caption{Cross-model semantic factorization. Seven classifiers with no syntactic overlap perform the same semantic operation on a shared parent. \tool synthesizes a new shared model that pre-aggregates all classifier flags in one pass.}
\label{fig:ex-factorization}
\end{figure}




The seven classifiers share no common subexpression, yet they all perform the same \emph{semantic operation}: for each order, decide whether any of its rows match a category-specific predicate.
\tool recognizes this cross-model pattern and factors it.
It synthesizes a new intermediate model \texttt{order\_flags} that, in a single pass over \texttt{claims}, computes one boolean flag per category by aggregating each category's predicate with \texttt{MAX(CASE WHEN ... THEN 1 ELSE 0 END)} grouped by \texttt{order\_id}.
Each original classifier is then rewritten as a lightweight lookup: select the orders whose corresponding flag in \texttt{order\_flags} is~$1$.
Seven full scans and seven independent grouping passes over \texttt{claims} are replaced with one pre-aggregation pass and seven inexpensive boolean filters.
 
This optimization requires three capabilities that go beyond traditional techniques.
First, the optimizer must recognize that seven syntactically dissimilar models perform the same semantic operation, which syntactic subexpression matching, the basis of multi-query optimization (MQO)~\cite{sellis1988multiple} and common subexpression elimination, cannot do.
Second, it must \emph{invent} a new intermediate model absent from the original project; materialized view recommendation tools~\cite{agrawal2000automated, bigsubs} synthesize candidate views, but they too rely on syntactic subexpression mining, of which the seven classifiers share none.
Third, it must rewrite the dependent models to consume the new intermediate while preserving end-to-end equivalence, a multi-model restructuring that single-query rewrite systems~\cite{wang2024wetune,zhou2024r} do not attempt.

%==============================================================
% EXAMPLE 2: Cross-Boundary Reasoning
%==============================================================
\head{Example 2: Cross-Boundary Semantic Reasoning}
The same project contains a classifier \texttt{nursing} that filters for skilled-nursing claims by facility code, then performs an anti-join against a sibling model \texttt{equipment} to exclude equipment-related claims (Figure~\ref{fig:ex-crossboundary}, left).
The anti-join forces the engine to materialize \texttt{equipment}, build a hash table, and probe it for every row of \texttt{claims}; it also introduces a DAG dependency that serializes the two models.



\begin{figure}[!htbp]
\begin{minipage}[t]{0.48\linewidth}
\figsubheader{Before: anti-join}
\begin{lstlisting}[language=SQL]
-- nursing
SELECT *
FROM {{ ref('claims') }} c
WHERE c.facility_code
        IN ('31','32')
  AND NOT EXISTS (
    SELECT 1
    FROM {{ ref('equipment') }} e
    WHERE e.order_id = c.order_id
      AND e.line_id  = c.line_id)
 
-- equipment (just a filter)
SELECT *
FROM {{ ref('claims') }}
WHERE category = 'equipment'
\end{lstlisting}
\end{minipage}
\hfill
\begin{minipage}[t]{0.48\linewidth}
\figsubheader{After: predicate}
\begin{lstlisting}[language=SQL]
-- nursing (rewritten)
SELECT *
FROM {{ ref('claims') }} c
WHERE c.facility_code
        IN ('31','32')
  AND (c.category IS NULL
    OR c.category <> 'equipment')


 
-- equipment is removed if no
-- other model depends on it;
-- otherwise its definition
-- is unchanged.
\end{lstlisting}
\vspace{1.6\baselineskip}
\end{minipage}
\caption{Cross-boundary semantic reasoning. By reading the SQL of the sibling \texttt{equipment} model, \tool determines that the anti-join is equivalent to a scalar predicate on a shared parent. The hash anti-join, the auxiliary scan, and the DAG dependency are all eliminated.}
\label{fig:ex-crossboundary}
\end{figure}




By reading the SQL definition of \texttt{equipment}, \tool discovers that it computes nothing more than a single-predicate filter on a shared parent.
The anti-join is therefore equivalent to an inline scalar predicate on \texttt{claims}, as shown on the right of the figure: the hash anti-join, the auxiliary scan, and the cross-model DAG dependency are all eliminated.
The same reasoning generalizes to every sibling classifier that uses \texttt{equipment} as an exclusion list, compounding the benefit across the pipeline.

The equivalence enabling this rewrite is not declared anywhere; it is \emph{embedded in the SQL of another model file}.
A traditional engine sees \texttt{equipment} either as a base table to be probed or as a view to be expanded into the local query, neither of which performs the cross-model analysis needed to eliminate the join entirely.
We acknowledge a tradeoff: the rewrite couples \texttt{nursing} to the current definition of \texttt{equipment}, so a future change to that definition would require updating the rewritten consumers.
In practice, stable upstream models such as \texttt{equipment} change rarely, and \tool re-evaluates the project after each developer change, so the cost of re-deriving the rewrite is amortized over many executions.

%==============================================================
% EXAMPLE 3: Frequency-Aware
%==============================================================
\head{Example 3: Frequency-Aware Restructuring}
Different sources in a dbt project update at different rates, and the structure of the DAG determines how much computation must be repeated at each high-frequency refresh.
Consider a model that joins fast-changing transactional data with slow-changing reference data and applies expensive classification logic to the slow-changing side (Figure~\ref{fig:ex-frequency}, left).


\begin{figure}[!htbp]
\begin{minipage}[t]{0.48\linewidth}
\figsubheader{Before: hourly classification}
\begin{lstlisting}[language=SQL]
-- order_summary  (hourly)
SELECT o.order_id, o.amount,
  CASE
    WHEN p.code IN (set_A) THEN 'A'
    WHEN p.code IN (set_B) THEN 'B'
    ...
  END AS category
FROM {{ ref('orders') }}  o
JOIN {{ ref('catalog') }} p
  ON o.product_id = p.product_id

  
-- orders:  hourly
-- catalog: weekly
\end{lstlisting}
\vspace{1.6\baselineskip}
\end{minipage}
\hfill
\begin{minipage}[t]{0.48\linewidth}
\figsubheader{After: split by frequency}
\begin{lstlisting}[language=SQL]
-- product_categories (weekly)
SELECT product_id,
  CASE
    WHEN code IN (set_A) THEN 'A'
    WHEN code IN (set_B) THEN 'B'
    ...
  END AS category
FROM {{ ref('catalog') }}
 
-- order_summary (hourly)
SELECT o.order_id, o.amount,
       pc.category
FROM {{ ref('orders') }} o
JOIN {{ ref('product_categories') }} pc
  USING (product_id)
\end{lstlisting}
\end{minipage}
\caption{Frequency-aware restructuring. Splitting the model along the frequency boundary moves expensive classification onto a weekly cadence; the hourly model becomes a lightweight join.}
\label{fig:ex-frequency}
\end{figure}


Because \texttt{orders} changes hourly, \texttt{order\_summary} re-executes hourly, and the multi-branch classification over \texttt{catalog} is recomputed on every hourly refresh, despite its inputs changing only weekly.
\tool traces the update frequency from each source through the DAG, identifies that this model has \emph{mixed-frequency inputs}, and splits it along the frequency boundary.
The expensive classification is extracted into a new slow-frequency model, and the original model becomes a lightweight equi-join.
When this pattern recurs across many consumers of the same slow-changing reference data, the savings compound: each hourly refresh avoids redundant recomputation of logic whose inputs have not changed.

This optimization is invisible to optimizers that treat cost as a fixed per-query quantity: the benefit of splitting the model is realized only when the two halves are known to execute at different rates.
Materialized view selection considers update cost but does not restructure the DAG to align computation with update frequency, and traditional query rewriters have no notion of refresh schedule at all.

\subsection{Why Existing Techniques Fall Short}
\label{sec:why-different}

The three examples above are not isolated curiosities; they illustrate a general capability gap.
Single-query rewrite systems~\cite{wang2024wetune,zhou2024r,chakravarthy1990logic} operate within the boundary of one query and cannot restructure across model files.
Multi-query optimization and common subexpression elimination~\cite{sellis1988multiple} share intermediate computations across queries but rely on syntactic matching, so they miss the cross-model semantic equivalence in Example~1.
Materialized view selection~\cite{agrawal2000automated} chooses among existing view candidates and treats update cost as a static parameter rather than as a structural driver of DAG design (Example~3).
Code clone detection and lineage tools can identify structural similarity or trace data provenance, but they do not perform the optimization itself.

The gap is not that any one of these techniques is poorly designed; it is that none of them targets the setting we are concerned with.
Optimizing a dbt project requires (1) cross-model semantic analysis, (2) DAG-aware structural rewriting that may invent or merge models, and (3) frequency-aware reasoning that ties structural decisions to update rates.
The remainder of this paper presents an approach that integrates these three capabilities into a coherent optimization framework.






\begin{comment}
    


\begin{figure}[t]
\figsubheader{Before: 7 independent classifiers (2 shown)}
% \centering\textbf{Before: 7 independent classifiers (2 shown)}
\begin{lstlisting}[language=SQL]
-- dialysis_classifier
SELECT DISTINCT order_id
FROM {{ ref('claims') }}
WHERE bill_code LIKE '72%'
   OR taxonomy IN (set_dialysis_tax)
   OR diagnosis IN (set_dialysis_dx)

-- psychiatric_classifier
SELECT DISTINCT order_id
FROM {{ ref('claims') }}
WHERE revenue IN (set_psych_rev)
   OR diagnosis BETWEEN '290' AND '319'
\end{lstlisting}

\figsubheader{After: shared pre-aggregation + lookups}
% \centering\textbf{After: shared pre-aggregation + lookups}
\begin{lstlisting}[language=SQL]
-- order_flags  (new shared model)
SELECT order_id,
  MAX(CASE WHEN bill_code LIKE '72%'
        OR taxonomy IN (set_dialysis_tax)
        OR diagnosis IN (set_dialysis_dx)
       THEN 1 ELSE 0 END) AS f_dialysis,
  MAX(CASE WHEN ... THEN 1 ELSE 0 END) AS f_psych,
  ...  -- flags for all 7 categories
FROM {{ ref('claims') }}
GROUP BY order_id

-- dialysis_classifier (rewritten)
SELECT order_id FROM {{ ref('order_flags') }}
WHERE f_dialysis = 1
\end{lstlisting}
\caption{Cross-model semantic factorization. Seven classifiers with no syntactic overlap perform the same semantic operation on a shared parent. \tool synthesizes a new shared model that pre-aggregates all classifier flags in one pass.}
\label{fig:ex-factorization}
\end{figure}

=========================




\begin{figure}[t]
\figsubheader{Before: anti-join against sibling model}
% \centering\textbf{Before: anti-join against sibling model}
\begin{lstlisting}[language=SQL]
-- nursing
SELECT *
FROM {{ ref('claims') }} c
WHERE c.facility_code IN ('31','32')
  AND NOT EXISTS (
    SELECT 1
    FROM {{ ref('equipment') }} e
    WHERE e.order_id = c.order_id
      AND e.line_id  = c.line_id)

-- equipment  (definition is just a filter)
SELECT * FROM {{ ref('claims') }}
WHERE category = 'equipment'
\end{lstlisting}

\figsubheader{After: anti-join replaced with predicate}
% \centering\textbf{After: anti-join replaced with predicate}
\begin{lstlisting}[language=SQL]
-- nursing  (rewritten)
SELECT *
FROM {{ ref('claims') }} c
WHERE c.facility_code IN ('31','32')
  AND (c.category IS NULL
    OR c.category <> 'equipment')

-- equipment is removed if no other model
-- depends on it; otherwise its definition
-- is unchanged.
\end{lstlisting}
\caption{Cross-boundary semantic reasoning. By reading the SQL of the sibling \texttt{equipment} model, \tool determines that the anti-join is equivalent to a scalar predicate on a shared parent. The hash anti-join, the auxiliary scan, and the DAG dependency are all eliminated.}
\label{fig:ex-crossboundary}
\end{figure}



=========================



\begin{figure}[t]
\figsubheader{Before: classification re-runs hourly}
% \centering\textbf{Before: classification re-runs hourly}
\begin{lstlisting}[language=SQL]
-- order_summary  (runs hourly)
SELECT o.order_id, o.amount,
  CASE
    WHEN p.code IN (set_A) THEN 'A'
    WHEN p.code IN (set_B) THEN 'B'
    ...
  END AS category
FROM {{ ref('orders') }}  o   -- hourly
JOIN {{ ref('catalog') }} p   -- weekly
  ON o.product_id = p.product_id
\end{lstlisting}

\figsubheader{After: split along frequency boundary}
% \centering\textbf{After: split along frequency boundary}
\begin{lstlisting}[language=SQL]
-- product_categories  (new, runs weekly)
SELECT product_id,
  CASE
    WHEN code IN (set_A) THEN 'A'
    WHEN code IN (set_B) THEN 'B'
    ...
  END AS category
FROM {{ ref('catalog') }}

-- order_summary  (rewritten, hourly, lightweight)
SELECT o.order_id, o.amount, pc.category
FROM {{ ref('orders') }} o
JOIN {{ ref('product_categories') }} pc
  USING (product_id)
\end{lstlisting}
\caption{Frequency-aware restructuring. Splitting the model along the frequency boundary moves expensive classification onto a weekly cadence; the hourly model becomes a lightweight join.}
\label{fig:ex-frequency}
\end{figure}









\end{comment}